\chapter{Conclusion}
In this thesis we have described an implementation of IETF Transport Services,
with a focus on how this can be used to expose the features of QUIC, TCP and
UDP through a common API. We have discussed the specific design considerations
which must be taken when implementing Transport Services in order to achieve as
flexible a system as possible. In this chapter we will discuss the findings
of the thesis overall and provide answers to the stated research questions.
In addition, we will discuss the scope and limitations of the thesis as a whole
and future work if CTaps is to be further developed into a production-ready library.

\section{Summary of Contributions}
The primary contribution of this thesis is CTaps.
CTas is an open source C-implementation of Transport Services as specified
by RFC~9621--9623. CTaps features support for QUIC, TCP and UDP as underlying
protocols and can therefore be used to analyse how well the features of these
protocols can be represented through a TAPS API. Internally, CTaps uses libuv
to serve as the central event loop allowing asyncronous interaction with the
network and Picoquic as a QUIC engine to provide support for QUIC.

CTaps implements large parts of the TAPS API as well
as the internal specification. Externally this means all major features of the
TAPS API except for peer-to-peer connection establishment. Internally this means
support for asynchronous candidate gathering, filtering, sorting and candidate racing. 
This means that a single CTaps client can connect to multiple servers supporting
different protocols. The effect of candidate racing is clearly illustrated
in \cref{fig:handshake-kde}.

In addition to the development of CTaps this thesis also contains
an extensive analysis of how QUIC maps onto TAPS and what architectural
challenges will have to be faced when implementing a TAPS system.
This is useful for potential future implementations of TAPS which
can use the lessons learned during the development of CTaps for their
own benefit.

\section{Research Questions and Findings}
At the onset of this thesis two primary research questions were defined
as goals and key motivations behind the development of CTaps. These
were:

\begin{enumerate}
    \item \rqone
    \item \rqtwo
\end{enumerate}

\subsection{Research Question One}
For the analysis of architecgture requirements and considerations
fo rTAPS implementations most of the requirements stem from
the need to be compliant with the architecture, functionality
and features as described in TAPS. The most important of these
are that TAPS is asynchronous at its core. This limits the actual
simplest naive implementation which is to simply interact with
the network in a blocking manner. As described in \cref{sec:async-background}
this means that TAPS implementations have two primary design
in order to achieve this asynchrony. Either use multiple threads
or be asynchronous.

The exact nature of which model should be folloewed
depends on the goals of the implementation. Asynchrony
is a more complicated model, both because of it deviation
from normal sequential programming, but also because
it requires an event loop \todo{always required?} which
requires a fairly heavy dependency.

Asynchronous models do however avoid the added overhead
of creating a large number of threads, as well as the
complexity needed for avoiding race conditions if the
implemnetation features any shared resources between the
threads.

Assuming that the implementation settles for an
asynchronous model like CTaps has then TAPS
requires that it pulls in an event loop
library such as libuv. 

Apart from the required asynchrony TAPS also mandates that the implementation
is flexible in which prtocol and andpoints a Connection can be attached to. Large parts
of the internal architcture is also prescribed through the candidate gathering and
candidate racing descriptions from RFC~9623. These prescribe how a Connection is created and how
the candidates are fetches. This is owever broad strokes and the exact implementation
is left up to the implementation. TAPS prescribes
nothing other than that it needs to be capable
of racing and able
to gather avrious candidates based
on the passed selection properties and
endpoints.

With the flexibility of binding between Conenction, Listener and underlying protocol and endpoints
required by CTaps there are also several implicit architectural constraints.
While not strictly required, in order to be as maintainable as possible
all TAPS implementaitons should feature some sort of generic way of representing
a protocol, such as the protocol interface in CTaps. This is to avoid the internal
implementation code becoming very tightly bound to specific protocols, making it
far more difficult to adapt the TAPS implementation to support new protocols.

In addition to the flexibility for established sockets, a true TAPS implementation
also needs a way to listen for incoming data for multiple different protocols
when opening a listener. Since RFC~9623 recommends that a listener should
be able to receive from multiple transport protocols at once over a single
socket \cite[Sec. 4.7]{rfc9623}, the internal structure used to represent protocols and manage
sockets need to reflect this advanced behavior. This means that the application
is barred from using simple abstractions such as the TCP and UDP handles in
libuv which do not support it. Instead it needs to implement more advanced
connection handling parsing incoming data directly.

With TAPS being connection-oriented it also mandates certain features
and therefore intermediary structures between the protocol and Connections.
In CTaps' case the requirement of TAPS being connection oriented comes into
play when supporing UDP. A TAPS implementation has to maintain its own Connections
for UDP, which can be done by emulating the four-tuple as used by TCP.
Incoming datagrams can be compared with a map of previous remote addresses
for the given local socket and demultiplexed accordingly. If no matching remote
can be found in the map then the datagram belongs to a new Connection.
This sets a hard requirement for internal data structures, as it
prescribes each socket having some sort of state in order to make these
lookups and that this state is stored together with the socket.
The requirements of multiple Connections sharing a single socket
also mandates handling which makes sure that the socket is not closed
until all associated Connections have been closed. This means that the
implementation will likely have to implement even more state in the form
of a counter or list keeping track of the number of open Connections still
using the given socket.

Since TAPS explicitly supports connection migration and multipath as well
as multiple local andpoints, the architecture must also be adopted to make
sure that the actual socket and Local Endpoint used for a given Connection
is flexible. This means that a Connection cannot be tightly bound to a single
endpoint if the TAPS implmentation wants to make use of connection migration
without featuring a workaround as in CTaps. Since TAPS also supports adding
endpoints during the lifetime of a Connection, it cannot rely on the Connection
always first binding to an ephemeral socket if it wants to make use of migration.

\subsection{Research Question Two}
While the TAPS API differs to a very large degree from the typical interface
used to interact with the transport layer which is that of a transport layer
protocol implementation. This means that there is certain control and
interactions which naturally differ between TAPS and actual transport protocols.
This is a very intentional feature of the design since TAPS is designed for the
application to interact with the network by stating its intentions instead of
requiring it to prescribe the exact behavior of the underlying protocol. This
is at the core of the second research question as there is naturally large
parts of the QUIC standard which is not directly controllable through TAPS,
but as this is intentional the more pertient analysis is whether the intentions
which can be registered with a TAPS implementation maps well to the features 
which are made availbale through QUIC, or if there are parts of the QUIC API
which is impossible ot access through the intention-based model of TAPS.

The primary result of this research question can be seen in \cref{tab:quic-taps-mapping}.
This table is a compilation of QUIC features as well as how they map to
the intention-based API of TAPS. The TAPS API being intention-based is
a key aspect here since there are QUIC features which are not avilable
through TAPS directly, but since TAPS is intention-based, the application
is able to state a desired result, and if the desired result can be expected
to be a result of the given QUIC feature, then a TAPS implementation can
reasonably make use of the given feature. A good example of this is the
settings related to ACKs, \mono{ack_delay_exponent} and \mono{max_ack_delay}.

While the application could in theory want to set the 
\mono{ack_delay_exponent} and \mono{max_ack_delay} transport parameters
for their own sake, it is far more likely it is far more likely that
these would be configured in order to acheive some downstream effect.
For the ACKs settings specifically, they could be modified in order
to for example optimize the latency of a connection, minimizing the
ack settings in order to be informed as quickly as possible about
a possible outage. TAPS reflects this not directly through the
exposing ACK settings directly but instead through the underlying
intention, allowing the appliaction to set the \mono{connCapacityProfile}
enumerator for the connection to Low Latency/Interactive. This will
cause the TAPS system to perform whatever it can to minimize the latency
for that given connection, including setting the ACK settings if available.

While this does not provide the granular control of setting ACK delays
explicitly, it also greatly simplifies the actual interaction and creates
a single unified interface for minimizing latency for all protocols.
The application does not have to perform research on which settings
actually reduce latency for the various protocols, instead having
it be handled by a single central library.

By looking through the this lens of intentionality and the desire
of the application, nearly all QUIC features can be made use of
by a TAPS system to various degrees. There are some concepts which
do not map cleanly such as starting a unidirectional stream on an
existing QUIC connection, but these are minor features attached
to a supported operation. If needed, these can be supported easily
by the TAPS implementation for example taking a flag parameters
which can be set to enable these specific featues.

When using TAPS and connecting to a peer using QUIC the expected
behavior of this conneciton is preserved well. For example multistreaming
is available through the clone API. This is the expected behavior
as laid out by the TAPS RFCs and it maps cleanly. Connection migration
is available thorugh using mulitple local endpoint and setting
the \mono{multipath} selection property to enabled. Even more granular
control is available thorugh the \mono{multipathPolicy}. These
controls ensure that the underlying QUIC connection does not
deviate from the behavior that is mandated by the application.

Since TPAS is an intention-based API then it naturally follows
that when using QUIC then QUIC too will act according to the
stated intentions of the application. This tracks well to the
insight gained from implementing TAPS with QUIC support
through CTaps. QUIC behaves as expected, with the single large
caveat of the connection being represented through the Connection Group
instead of the Connection object.

When it comes to the internal architecture and mapping of objects
there are also certain parts of QUIC which a TAPS implementation
has to take into account if it wants to expose as many features
of QUIC as possible. The largest of these features is multistreaming,
which requires that the underlying QUIC connection is treated as a
Connection Group in TAPS, since this is the abstraction TAPS
uses for connected streams. As is dicussed in \cref{sec:quic-streams-under-taps}
this means that a decision has to be made surrounding the semantics for closing
a TAPS Connection when QUIC is the underlying protocol. If the
implementation desired to present exactly the same API at all
points with no deviations then each stream must be forcefully closed
through a \mono{STOP_SENDING}, in order to have \mono{close()}
close the Connection in both directions, as is expected for TCP
and UDP.

If the implementation instead, like CTaps, prioritizes having
\mono{close()} represent only a graceful teardown, then \mono{close()}
has to be reinterpreted to mean that the application is done
writing on a Connection if QUIC is the underlying connection.
Since an application can achieve unified behavior through always
interacting with a Connection Group then this is not a breaking
issue but it is a deviation from how a application would expect
to interact with the network coming purely from TCP.

An additional consideratoin when adding support for QUIC to a
TAPS implementation is that QUIC connctions always transfer
data over streams. This means that if the implementation
wishes to consistenly map a QUIC connection to a connection group
and a stream to a Connection, then all Connections will have
to belong to a Connection Group if they are based on top of
QUIC. CTaps has chosen to follow the pattern of a Connection
always belonging to a Connection Group for all Connections, no
matter what the underlying protocol is. This greatly simplifies
internal logic related to Connection Groups. This also greatly
simplifies the implementation of the single unified way of
interacting with Connection Groups, which ensures that applications
can be written without taking the underlying protocol into account.

\section{Scope and Limitations}
CTaps is an exploratory work, intended to show how an implementation of a TAPS
system might look in a real world context, but is not intended to be used
in a production context.

CTaps has throughout had a larger focus on exposing as much as possible of
the TAPS API, without necessarily implementing the internal logic associated with
the actual API settings. 
This was done to see exactly how the designed TAPS
RFC~would actually map into code and that the different settings could actually be
presented well. Although this has proven successful in providing a public
API to reason about and guidance for future TAPS implementations, this is
not the ideal starting location when considering future real-world usage of
CTaps.

Application developers coming across CTaps would likely find the API to be fairly confusing.
It looks like a more complete implementation than it actually is, and several
attempts to configure connections with unsupported settings, for example the multipath
policy connection property, would lead to no changed behavior from the applications point of view.

Another consequence of this push to make CTaps present as much of the API of TAPS
as possible is that most testing and development has been focused on the various
happy paths of the features, while less time or testing has been devoted
to testing various error paths, both for errors stemming from errors in user configuration
and from network errors.

Another key limitation with the current implementation of CTaps is that several key
features have been deemed to be out of scope, and are not presented in the API at
all. The largest group of these is peer to peer connection establishment. This is
the only major group of TAPS features which is totally absent from CTaps.

Among the smaller features not present in CTaps, prominent examples are that CTaps
does not support adding endpoints in during the lifetime of a connection and that
CTaps does not feature
listening on more than a single connection.

Not being able to add endpoints to an active Connection places a real limitation on
a Connection's ability to migrate, as there is no way to dynamically administer the
endpoints which can be migrated to or from. All Local and Remote endpoints have to be
passed to the Preconnection. QUIC connections can still migrate to new remote endpoints,
if migration is initiated by the peer, but there is no way to dynamically
add a new endpoint for example.

CTaps also contain inaccuracies which deviate from expected behavior, for example
that \mono{sent()}-callbacks
do not represent when a packet has been actually transferred to the network,
but instead when a packet has been handed over to Picoquic if QUIC is the underlying
protocol. QUIC may then cache the packets internally and the application
cannot therefore use the sent callbacks to reliably get the state of the send buffers.

\section{Future Work}\label{sec:future-work}
If an attempt was to be made to make CTaps into a real-world production
ready system, then candidate interleaving would likely be one of the first
features to be implemented.

It would have to be carefully considered exactly how these candidates were
to be interleaved. The applications does rank its most desired protocols,
so interleaving with no care as to the order could lead to a lower ranked
protocol being tried second after the most desirable one, to the application's
surprise. It would also be beneficial to include the Local and Remote Endpoints
in the interleaving algorithm somehow. In the hypothetical situation that the
first candidates all try to connect to the same Eemote Endpoint, but that endpoint
turns out to be unreachable, the current stagger delay would lead to a one
second delay before the next endpoint was tried.

In order to mitigate the possible DNS delay,
connection attempts in CTaps can be made the same way as recommended in Happy Eyeballs V2, racing
as soon as a resolved endpoint comes back, with a \qty{50}{\milli\second} resolution delay
to give preferred addresses a chance to return before initiating an
attempt with a less desirable endpoint.

\todo{Describe more flexible socket manager handling}

If CTaps were to be developed into a fully featured TAPS library, it would likely
be wise to first strip away the parts of the external API which currently have
no support in the internal implementation. This would allow a gradual addition of
features as they were supported, making it more clear to downstream applications exactly
which parts of the external library is mature enough to be used. This would also allow
future development to focus on more important improvements, instead of implementing
as much of the API as possible in order to make the features exposed by the API actually
real. The highest priority future work is described in \cref{sec:future-work}.

\subsection{Alternative QUIC Backends} \todo{This section is valuable but needs to be tightened by a lot}
Since QUIC is an exploratory work an prioritized simplicity over raw performance
when selecting the QUIC library to use there are likely large performance gains
to be made by selecting a more performance oriented library than Picoquic. While
CTaps maintains a generic interface to the underlying protocols and the switch
would therefore in theory consist of only rewriting the two QUIC files themselves,
This requires that the QUIC library would be compatible with the "QUIC engine"
mode of operation as epected by CTaps, in order to integrate it into the libuv
event loop.

An alternative more performance-oriented QUIC library in C is for example
be MsQuic~\cite{MsQuic}. MsQuic is a cross-platform C implementation of QUIC
developed by Microsoft. With MsQuic being a cross-platform library it would
further enable CTaps to become cross-platform as well, since the core dependencies
of Libuv and MsQuic would be cross-platform. For CTaps, the QUIC
implementation of the protocol interface is already by far the most complicated
of the current protocol implementations as seen in \cref{tab:loc_complexity}.
MsQuic is a far larger library than Picoquic at $\sim135~000$ LoC compared to
Picoquic's $\sim84~000$. While this means that more features and performance
is available, this would also likely lead to a larger amount of complexity
in the configuration from CTaps' point of view, inflating the difference
in the complexity in the various protocol implementations even further.

In normal execution mode MsQuic owns the event loop, notifying
the application of updates and asking for input through registered
callbacks. This mirrors how CTaps interacts with its downstream
applications. In theory this execution mode could be integrated
into CTaps, but it would require using async handles in libuv [cite], since
MsQuic specifically require that the callback functions do not
introduce significant delay~\cite{msquic_execution}. This would
be a significant architectural change in CTaps and introduce
a large amount of complexity.

MsQuic also features an alternative execution mode where the
application, in this case CTaps, creates the threads it will
run on and periodically call into MsQuic to perform thhe necessary
actions. This mirrors the timed wakeups required by Picoquic
and could be integrated with the existing functionality in
CTaps. The largest difference with the current implementation
is that CTaps would use create a event queue object
through picoquic instead of using a timer to wake up at set times.
CTaps would then have to monitor this event queue object with \mono{uv_poll_t}
instead of using a \mono{uv_timer_t} handle as with Picoquic.
Examples from the MsQuic documentation creating a event
queue object and draining it can be seen in \cref{lst:msquicCustomEpoll}
and \cref{lst:msquicCustomLoop}.

While this shows that CTaps is not very tightly bound
to Picoquic itself, it is very tightly bound to libuv and
requires any additional libraries to be able to integrate
well with an existing event loop. For example if 
MsQuic did not feature the alternative execution mode
then integrating it with CTaps would be a large undertaking,
likely changing the entire internal architecture to
make the updates form QUIC flow through libuv async
handles.

\begin{listing}[htbp]
  \begin{minted}{C}
HANDLE IOCP = CreateIoCompletionPort(INVALID_HANDLE_VALUE, nullptr, 0, 1);
QUIC_EXECUTION_CONFIG ExecConfig = { 0, &IOCP };

QUIC_EXECUTION* ExecContext = nullptr;
QUIC_STATUS Status = MsQuic->ExecutionCreate(QUIC_GLOBAL_EXECUTION_CONFIG_FLAG_NONE, 0, 1, &ExecConfig, &ExecContext);
\end{minted}
\caption{MsQuic code to create an event queue object~\cite{msquic_execution}.}
\label{lst:msquicCustomEpoll}
\end{listing}

\begin{listing}[htbp]
  \begin{minted}{C}
bool AllDone = false;
while (!AllDone) {
    uint32_t WaitTime = MsQuic->ExecutionPoll(ExecContext);

    ULONG OverlappedCount = 0;
    OVERLAPPED_ENTRY Overlapped[8];
    if (GetQueuedCompletionStatusEx(IOCP, Overlapped, ARRAYSIZE(Overlapped), &OverlappedCount, WaitTime, FALSE)) {
        for (ULONG i = 0; i < OverlappedCount; ++i) {
            QUIC_SQE* Sqe = CONTAINING_RECORD(Overlapped[i].lpOverlapped, QUIC_SQE, Overlapped);
            Sqe->Completion(&Overlapped[i]);
        }
    }
}
\end{minted}
\caption{MsQuic code to poll custom event queue object~\cite{msquic_execution}.}
\label{lst:msquicCustomLoop}
\end{listing}
  

